% Solutions of 19-sorting.
\section*{Chapter \ref{ch:b1:sorting} --- Sorting}

\begin{solution}{exo:b1:sorting:1}
The \omterm{def:b1:sorting:quicksort}{pivot} is 6. The keys 3, 2, 5, 1, 4 are smaller and are moved, in this order, to the front; 8 and 7 are not. The
final exchange puts 6 after them: $[3, 2, 5, 1, 4, 6, 7, 8]$, \omterm{def:b1:sorting:quicksort}{pivot} at index 5, after 7 comparisons.
\end{solution}

\begin{solution}{exo:b1:sorting:2}
1 (1 against 2), 2 (4 against 2), 3 (4 against 3), 4 (4 against 8), 7 (7 against 8), 8 (9 against 8), 9 (9 against 10);
the first list is then empty and 10, 11 are copied: $[1, 2, 3, 4, 7, 8, 9, 10, 11]$ in 7 comparisons.
\end{solution}

\begin{solution}{exo:b1:sorting:3}
Counts for 0 to 5: $1, 2, 1, 2, 1, 1$. Block starts: $0, 1, 3, 4, 6, 7$. Output: $[0, 1, 1, 2, 3, 3, 4, 5]$, the two 1s
and the two 3s in their input order.
\end{solution}

\begin{solution}{exo:b1:sorting:4}
Only $(1, \mathrm{b}), (1, \mathrm{d}), (2, \mathrm{a}), (2, \mathrm{c})$. Selection sort is not \omterm{def:b1:sorting:stable}{stable}: on $(2,
\mathrm{a}), (2, \mathrm{c}), (1, \mathrm{b})$, the first exchange moves $(2, \mathrm{a})$ behind $(2, \mathrm{c})$, and
the output is $(1, \mathrm{b}), (2, \mathrm{c}), (2, \mathrm{a})$.
\end{solution}

\begin{solution}{exo:b1:sorting:5}
With all keys equal, no key is smaller than the \omterm{def:b1:sorting:quicksort}{pivot}: the partition makes $m - 1$ comparisons, leaves the left part
empty and the right part with $m - 1$ keys. In all $(n - 1) + \dots + 1 = n(n - 1)/2$, which is 499\,500 for $n =
1000$. The \omterm{def:b1:sorting:quicksort}{three-way partition} compares each key twice (not smaller, not larger), puts all of them in the middle, and
recurses on two empty parts: $2n$.
\end{solution}

\begin{solution}{exo:b1:sorting:6}
Each key goes to the last free cell of its digit's block, so the blocks are in digit order: one pass sorts by its digit.
But within a block the keys come out in the reverse of their input order: the sort is not \omterm{def:b1:sorting:stable}{stable}. The last pass, on the
hundreds, reverses 329 and 355, which the pass on the tens had put in order; the result is wrong. The repair walks the
input from its end, \texttt{for (int i = n - 1; i >= 0; i-{}-)}, so that the last key of each digit takes the last cell.
\end{solution}

\begin{solution}{exo:b1:sorting:7}
Sort by date first, then by member, with a \omterm{def:b1:sorting:stable}{stable sort}: the second sort orders by member, and stability keeps the date
order inside each member (\code{exercises/loans_by_member.py}). In the other order the dates would be the primary key.
The test compares the result with a sort on the pair (member, date).
\end{solution}

\begin{solution}{exo:b1:sorting:8}
Exchanging two neighbours $a_j > a_{j+1}$ removes the \omterm{def:b1:sorting:inversion}{inversion} they form and changes no other pair's relative order:
the number of \omterm{def:b1:sorting:inversion}{inversions} drops by exactly 1. Insertion sort exchanges only such neighbours, and stops with a sorted
\omterm{def:b1:pointers-and-arrays:array}{array}, which has no \omterm{def:b1:sorting:inversion}{inversion}. So it makes exactly as many exchanges as there were \omterm{def:b1:sorting:inversion}{inversions}.
\end{solution}

\begin{solution}{exo:b1:sorting:9}
$12! = 479\,001\,600$ and $\log_2 12! \approx 28.84$: every \omterm{def:b1:sorting:comparison-sort}{comparison sort} needs at least 29 comparisons for some order
of 12 keys. \omterm{def:b1:sorting:merge-sort}{Merge sort} needs 33 in its worst case, \omterm{def:b1:sorting:merge-sort}{merge} insertion 30. The bound does not say whether 29 can be reached.
\end{solution}

\begin{solution}{exo:b1:sorting:10}
Multiplying by $n$: $nC(n) = n(n - 1) + 2\sum_{k<n} C(k)$; the same for $n - 1$, subtracted:
$nC(n) - (n - 1)C(n - 1) = 2(n - 1) + 2C(n - 1)$, so $nC(n) = (n + 1)C(n - 1) + 2(n - 1)$. Dividing by $n(n + 1)$ and
using $\frac{2(n - 1)}{n(n + 1)} = \frac{4}{n + 1} - \frac2n$:
\[ \frac{C(n)}{n + 1} = \frac{C(n - 1)}{n} + \frac{4}{n + 1} - \frac2n = \sum_{j=2}^{n} \Bigl(\frac{4}{j + 1} -
\frac2j\Bigr) = 2H_n + \frac{4}{n + 1} - 4, \]
since $C(1) = 0$. Hence $C(n) = 2(n + 1)H_n - 4n$; with $H_n \approx \ln n$, $C(n) \approx 2n \ln n$.
\end{solution}

\begin{solution}{exo:b1:sorting:11}
\omcode{code/bachelor-1/19-sorting/exercises/inversions.c}{5}{14}{Counting \omterm{def:b1:sorting:inversion}{inversions} while merging}
An \omterm{def:b1:sorting:inversion}{inversion} has both ends in the left half, both in the right half, or one in each; the recursive calls count the first
two kinds. When the \omterm{def:b1:sorting:merge-sort}{merge} takes $a[j]$ from the right half while $h - i$ keys remain in the left half, each of them is
larger than $a[j]$ and before it: $h - i$ \omterm{def:b1:sorting:inversion}{inversions} of the third kind, each counted once. The cost is that of \omterm{def:b1:sorting:merge-sort}{merge
sort}, $O(n \log n)$; the test compares with the double loop for $n \leq 300$ on six kinds of input.
\end{solution}

\begin{solution}{exo:b1:sorting:12}
For $j = 0$ there is nothing to prove. Suppose the keys sorted by their last $j$ digits before the pass on digit $j + 1$.
Take two keys $x$ and $y$. If their digits $j + 1$ differ, the pass orders them by that digit, which decides the order of
their last $j + 1$ digits. If they are equal, the pass, being \omterm{def:b1:sorting:stable}{stable}, leaves $x$ and $y$ in their previous order, the
order of their last $j$ digits, which is then the order of their last $j + 1$ digits. So after the pass the keys are
sorted by their last $j + 1$ digits.
\end{solution}

\begin{solution}{pb:b1:sorting:1}
\textbf{1.} $2^6 = 64 < 120 \leq 128 = 2^7$: $\lceil \log_2 120 \rceil = 7$. A tree of height 6 has at most 64 leaves,
fewer than the 120 orders.

\textbf{2.} The height of the tree is the worst case. The average depth of 120 leaves is at least $\log_2 120 \approx
6.91$; \omterm{def:b1:sorting:merge-sort}{merge} insertion's mean, 6.93, is almost there.

\textbf{3.} At least 120 leaves; at most $2^7 = 128$.

\textbf{4.} A \omterm{def:b1:induction-and-termination:rooted-tree}{leaf} at depth $d \leq 6$ takes at least 2 of the 128 positions at depth 7, a \omterm{def:b1:induction-and-termination:rooted-tree}{leaf} at depth 7 one: $2x +
(120 - x) \leq 128$, so $x \leq 8$. \omterm{def:b1:sorting:merge-sort}{Merge} insertion has exactly 8 orders sorted in 6 comparisons.

\textbf{5.} $a < b < d$ and $c < d$ leave 3 orders of $a, b, c, d$ ($c$ before $a$, between $a$ and $b$, or between $b$ and
$d$), and $e$ can be in any of 5 places: 15 orders. After three comparisons, the leaves still reachable must fit in a
\omterm{def:b1:induction-and-termination:rooted-tree}{subtree} of height 4, which has at most 16 leaves.

\textbf{6.} $\lceil \log_2 24 \rceil = 5$. \omterm{def:b1:sorting:merge-sort}{Merge sort} sorts 4 keys in at most $1 + 1 + 3 = 5$ comparisons.

\textbf{7.} 10 on the reversed order $(4, 3, 2, 1, 0)$, which has 10 \omterm{def:b1:sorting:inversion}{inversions}; 4 on the sorted order.

\textbf{8.} Insertion sort makes one comparison per exchange, plus one per element that stops before reaching the
front. On average there are $n(n - 1)/4 = 5$ \omterm{def:b1:sorting:inversion}{inversions}, and the element at position $i$ reaches the front with
probability $1/(i + 1)$: $5 + 4 - (\frac12 + \frac13 + \frac14 + \frac15) \approx 7.72$.

\textbf{9.} $W(1) = 0$, $W(2) = 1$, $W(3) = W(1) + W(2) + 2 = 3$, $W(5) = W(2) + W(3) + 4 = 8$.

\textbf{10.} The order $(0, 3, 2, 1, 4)$ (found by the program): every \omterm{def:b1:sorting:merge-sort}{merge} runs to its last element.

\textbf{11.} Each \omterm{def:b1:sorting:merge-sort}{merge} of $p$ and $q$ elements makes at least $\min(p, q)$ comparisons: $1 + (1 + 1) + 2 = 5$.

\textbf{12.} \omterm{def:b1:sorting:merge-sort}{Merge sort}: 8 against 10, one above the bound.

\textbf{13.} $a < b < d$ and $c < d$; nothing about $e$.

\textbf{14.} Binary search in a chain of 3 has 4 possible places, and $\lceil \log_2 4 \rceil = 2$.

\textbf{15.} $c < d$, so its place is among the keys below $d$: $a$, $b$, and $e$ if $e < d$. Three keys give 4 places,
2 comparisons.

\textbf{16.} $3 + 2 + 2 = 7$.

\textbf{17.} $a, b, c, d, e = 4, 1, 3, 0, 2$. Comparison 1 exchanges $a$ and $b$ ($a = 1$, $b = 4$); comparison 2
exchanges $c$ and $d$ ($c = 0$, $d = 3$); comparison 3, $3 < 4$, exchanges the pairs: $a, b, c, d = 0, 3, 1, 4$. The
chain $0 < 3 < 4$; $e = 2$ is compared with 3, then 0: chain $0, 2, 3, 4$. Then $c = 1$, below 4, is compared with 2,
then 0: $0, 1, 2, 3, 4$, in 7 comparisons.

\textbf{18.} Inserting $c$ first costs 2 comparisons (among $a$, $b$), but leaves a chain of 4 for $e$, with 5 places:
3 comparisons, 8 in all. Inserting $e$ first keeps both insertions at 4 places.

\textbf{19.} That the result is sorted, and how many comparisons it took.

\textbf{20.} A \omterm{def:b1:sorting:comparison-sort}{comparison sort}'s behaviour depends only on the outcomes of its comparisons, that is on the relative order
of the keys; 120 orders are all the relative orders of five distinct keys.

\textbf{21.} 8 orders. Yes: the bound is on the worst case, and question 4 allows up to 8.

\textbf{22.} The terms for $k = 1, \dots, 5$ are $0, 1, 2, 2, 2$: 7. For 12 keys the method needs 30, one more than the
bound 29.

\textbf{23.} $20! \approx 2.4 \times 10^{18}$ orders, far too many to enumerate.

\textbf{24.} \emph{Five keys can be sorted in seven comparisons, and no method does it in six}: six comparisons
distinguish at most 64 outcomes, fewer than the 120 orders (question 1); \omterm{def:b1:sorting:merge-sort}{merge} insertion sorts every one of the 120
orders in at most seven, as the exhaustive check confirms (\cref{out:b1:sorting:sort5}).
\end{solution}
